\documentclass[journal,12pt,draftclsnofoot,onecolumn]{IEEEtran}

\usepackage{graphicx}
\usepackage{booktabs}
\usepackage{float}

\begin{document}

\title{Minimal IEEE One-Column Draft}

\author{
    \IEEEauthorblockN{First Author\textsuperscript{1}, Second Author\textsuperscript{2}}
    \IEEEauthorblockA{\textsuperscript{1}Department of Examples, Example University\\
    \textsuperscript{2}Department of Conversion, Sample Hospital}
}

\maketitle

\begin{abstract}
This minimal draft intentionally shows common one-column IEEEtran patterns that need attention before a two-column journal submission.
\end{abstract}

\begin{IEEEkeywords}
IEEEtran, one-column draft, two-column conversion.
\end{IEEEkeywords}

\section{Introduction}
One-column drafts are convenient during writing, but some draft patterns do not transfer cleanly to IEEE journal two-column output.

\begin{figure}[H]
\centering
\fbox{\rule{0pt}{1.2in}\rule{0.92\textwidth}{0pt}}
\caption{A draft figure pinned with \texttt{[H]} and sized to \texttt{\textbackslash textwidth}.}
\label{fig:draft}
\end{figure}

\begin{table}[H]
\centering
\caption{A draft table pinned with \texttt{[H]}.}
\label{tab:draft}
\begin{tabular}{lcc}
\toprule
Method & Dice & HD95 \\
\midrule
Baseline & 0.84 & 41.2 \\
Converted & 0.86 & 40.5 \\
\bottomrule
\end{tabular}
\end{table}

\section{Conclusion}
This file is intentionally shaped like a draft that should be converted before submission.

\begin{thebibliography}{1}
\bibitem{ieeetran}
M.~Shell, ``IEEEtran,'' CTAN, 2015.
\end{thebibliography}

\end{document}
